% 本文件由 paper/generate-content.js 从游戏数据自动生成，请勿手工编辑。

\clearpage
\subsection{赵敏线：朔漠玫瑰}

\subsubsection{第一拍·绿柳初逢}\label{plot:zhaomin:1}
本关以选择判断完成叙事向健康知识的转场，题目见\hyperref[question:zhaomin:1]{第\ref*{question:zhaomin:1}项}。

\plotframe{assets/levels/zhaomin/level_01/frame_01.jpg}
  {赵敏线第一拍《绿柳初逢》第1帧}
  {fig:level:zhaomin:1:1}
  {\scriptsize}
  {\textbf{旁白：}元顺帝至正年间，天下大乱，群雄并起。蒙古汝阳王府郡主赵敏，奉旨南下，欲收服中原武林。这一日，她在绿柳山庄设下机关，引明教教主张无忌入局……\par}

\plotframe{assets/levels/zhaomin/level_01/frame_02.jpg}
  {赵敏线第一拍《绿柳初逢》第2帧}
  {fig:level:zhaomin:1:2}
  {\scriptsize}
  {\textbf{赵敏：}张教主，这地牢之中，你我不如聊些别的。你可知——人的足底，藏着全身的奥秘？\par
\textbf{旁白：}赵敏所问，正是中医学中极为重要的足底穴位——涌泉穴。此穴位于足底前三分之一凹陷处，乃是肾经首穴。传统医籍常以其补肾宁神、引火归元，女子日常按揉可作保健辅助，若有月经、睡眠或腰膝等不适，仍应辨证求医。\par}

\plotframe{assets/levels/zhaomin/level_01/frame_03.jpg}
  {赵敏线第一拍《绿柳初逢》第3帧}
  {fig:level:zhaomin:1:3}
  {\scriptsize}
  {\textbf{张无忌：}郡主这是何意？\par
\textbf{赵敏：}我蒙古萨满常说，足底有穴，通五脏六腑。方才你碰到我这里——（指了指足心）可觉得是什么穴？\par}

\clearpage

\plotframe{assets/levels/zhaomin/level_01/frame_04.jpg}
  {赵敏线第一拍《绿柳初逢》第4帧}
  {fig:level:zhaomin:1:4}
  {\scriptsize}
  {\mbox{}}

\plotframe{assets/levels/zhaomin/level_01/frame_05.jpg}
  {赵敏线第一拍《绿柳初逢》第5帧}
  {fig:level:zhaomin:1:5}
  {\scriptsize}
  {\mbox{}}

\subsubsection{第二拍·黑玉断续}\label{plot:zhaomin:2}
本关以选择判断完成叙事向健康知识的转场，题目见\hyperref[question:zhaomin:2]{第\ref*{question:zhaomin:2}项}。

\plotframe{assets/levels/zhaomin/level_02/frame_01.jpg}
  {赵敏线第二拍《黑玉断续》第1帧}
  {fig:level:zhaomin:2:1}
  {\scriptsize}
  {\textbf{旁白：}武当山上，俞岱岩师叔因二十年前被大力金刚指捏碎全身骨骼，瘫痪至今。赵敏为表诚意，携西域奇药黑玉断续膏而来。\par
\textbf{赵敏：}张真人，此乃黑玉断续膏，以西域秘法炼制，可接续断骨。但敷用之时，需以深厚内功催动药力入骨，辅以正骨手法。\par}

\clearpage

\plotframe{assets/levels/zhaomin/level_02/frame_02.jpg}
  {赵敏线第二拍《黑玉断续》第2帧}
  {fig:level:zhaomin:2:2}
  {\scriptsize}
  {\textbf{张三丰：}老道行医数十载，知接骨之要，在于骨正筋柔，气血以流。郡主此药，可含自然铜、骨碎补之类？\par
\textbf{赵敏：}真人明鉴。我虽不懂医道，但也听王府医官说过——女子年过七七（四十九岁），天癸竭，地道不通。以今日医理观之，绝经后骨量流失会加快；俞三侠虽非年迈，但卧床二十年，骨量下降风险甚高，需格外小心才是。\par}

\plotframe{assets/levels/zhaomin/level_02/frame_03.jpg}
  {赵敏线第二拍《黑玉断续》第3帧}
  {fig:level:zhaomin:2:3}
  {\scriptsize}
  {\textbf{旁白：}赵敏提及的骨质疏松，正是当代女性健康的重要议题。中医以肾主骨为理论核心，认为肾精充足则骨骼强健；西医则强调雌激素对骨密度的保护作用。绝经后女性因雌激素水平下降，骨质疏松风险显著上升。\par}

\plotframe{assets/levels/zhaomin/level_02/frame_04.jpg}
  {赵敏线第二拍《黑玉断续》第4帧}
  {fig:level:zhaomin:2:4}
  {\scriptsize}
  {\mbox{}}

\clearpage

\plotframe{assets/levels/zhaomin/level_02/frame_05.jpg}
  {赵敏线第二拍《黑玉断续》第5帧}
  {fig:level:zhaomin:2:5}
  {\scriptsize}
  {\mbox{}}

\subsubsection{第三拍·万安寺火}\label{plot:zhaomin:3}
本关以选择判断完成叙事向健康知识的转场，题目见\hyperref[question:zhaomin:3]{第\ref*{question:zhaomin:3}项}。

\plotframe{assets/levels/zhaomin/level_03/frame_01.jpg}
  {赵敏线第三拍《万安寺火》第1帧}
  {fig:level:zhaomin:3:1}
  {\scriptsize}
  {\textbf{旁白：}万安寺大火，六大派高手被困塔顶。火焰自下而上燃烧，浓烟弥漫。这场大火不仅烧毁了一座名刹，更让赵敏第一次直面生死。\par
\textbf{赵敏：}（望着火光，喃喃自语）父王要我收服武林……可若人都烧死了，收服又有何用？\par}

\plotframe{assets/levels/zhaomin/level_03/frame_02.jpg}
  {赵敏线第三拍《万安寺火》第2帧}
  {fig:level:zhaomin:3:2}
  {\scriptsize}
  {\textbf{旁白：}大火之中，多人被烧伤。烧伤急救——这一在现代医学中至关重要的课题，在这元末乱世之中，却少有人精通。\par
\textbf{赵敏：}张无忌！接着——（抛出一个药囊）这是王府的金创药，对烧伤有奇效！记住——先用清水冲洗伤口，再敷此药！\par}

\clearpage

\plotframe{assets/levels/zhaomin/level_03/frame_03.jpg}
  {赵敏线第三拍《万安寺火》第3帧}
  {fig:level:zhaomin:3:3}
  {\scriptsize}
  {\mbox{}}

\plotframe{assets/levels/zhaomin/level_03/frame_04.jpg}
  {赵敏线第三拍《万安寺火》第4帧}
  {fig:level:zhaomin:3:4}
  {\scriptsize}
  {\mbox{}}

\plotframe{assets/levels/zhaomin/level_03/frame_05.jpg}
  {赵敏线第三拍《万安寺火》第5帧}
  {fig:level:zhaomin:3:5}
  {\scriptsize}
  {\mbox{}}

\clearpage

\subsubsection{第四拍·灵蛇海上}\label{plot:zhaomin:4}
本关以多轮简答完成叙事向健康知识的转场，题目见\hyperref[question:zhaomin:4]{第\ref*{question:zhaomin:4}项}。

\plotframe{assets/levels/zhaomin/level_04/frame_01.jpg}
  {赵敏线第四拍《灵蛇海上》第1帧}
  {fig:level:zhaomin:4:1}
  {\scriptsize}
  {\textbf{旁白：}灵蛇岛之行，海路迢迢。赵敏自幼生长于蒙古草原，不惯乘船。海上颠簸数日，终感身体不适……\par}

\plotframe{assets/levels/zhaomin/level_04/frame_02.jpg}
  {赵敏线第四拍《灵蛇海上》第2帧}
  {fig:level:zhaomin:4:2}
  {\scriptsize}
  {\textbf{小昭：}赵姑娘，你这几日食欲不振，又常恶心呕吐……莫非是——\par
\textbf{赵敏：}（面色苍白，勉强一笑）莫要瞎说。我只是晕船罢了。小昭，你可知道治晕船的法子？\par
\textbf{小昭：}波斯医者常以生姜、薄荷止呕。我煮了姜汤，你且喝下。不过……赵姑娘，你当真只是晕船么？女子若有身孕，早期症状与晕船颇为相似——\par}

\plotframe{assets/levels/zhaomin/level_04/frame_03.jpg}
  {赵敏线第四拍《灵蛇海上》第3帧}
  {fig:level:zhaomin:4:3}
  {\scriptsize}
  {\mbox{}}

\clearpage

\plotframe{assets/levels/zhaomin/level_04/frame_04.jpg}
  {赵敏线第四拍《灵蛇海上》第4帧}
  {fig:level:zhaomin:4:4}
  {\scriptsize}
  {\mbox{}}

\plotframe{assets/levels/zhaomin/level_04/frame_05.jpg}
  {赵敏线第四拍《灵蛇海上》第5帧}
  {fig:level:zhaomin:4:5}
  {\scriptsize}
  {\mbox{}}

\subsubsection{第五拍·弥勒佛庙}\label{plot:zhaomin:5}
本关以选择判断完成叙事向健康知识的转场，题目见\hyperref[question:zhaomin:5]{第\ref*{question:zhaomin:5}项}。

\plotframe{assets/levels/zhaomin/level_05/frame_01.jpg}
  {赵敏线第五拍《弥勒佛庙》第1帧}
  {fig:level:zhaomin:5:1}
  {\scriptsize}
  {\textbf{旁白：}离开灵蛇岛后，赵敏为救张无忌，与父兄决裂。二人在弥勒佛庙中躲避追兵。张无忌身受刀伤，失血不止……\par}

\clearpage

\plotframe{assets/levels/zhaomin/level_05/frame_02.jpg}
  {赵敏线第五拍《弥勒佛庙》第2帧}
  {fig:level:zhaomin:5:2}
  {\scriptsize}
  {\textbf{赵敏：}无忌，你这伤口颇深，需先止血。我随身带有金创药——但需先清洗伤口，确认无异物残留，方可敷药。你可记得蝴蝶谷胡青牛教你的外伤治法？\par
\textbf{张无忌：}（虚弱地）外伤处理……先以清水或酒冲洗……若伤口深，需缝合……止血可用按压法，或在伤口近心端加压……金创药中，以三七、白及止血最佳……\par}

\plotframe{assets/levels/zhaomin/level_05/frame_03.jpg}
  {赵敏线第五拍《弥勒佛庙》第3帧}
  {fig:level:zhaomin:5:3}
  {\scriptsize}
  {\textbf{赵敏：}好。我先以洁净清水洗去污物，再用干净布巾持续按压止血；若伤口深、污染重，仍须尽快寻医。只是——你可知女子月经过多，亦会造成失血？我蒙古女子常谈当归、艾叶之类，你们中原医家又如何辨证？\par}

\plotframe{assets/levels/zhaomin/level_05/frame_04.jpg}
  {赵敏线第五拍《弥勒佛庙》第4帧}
  {fig:level:zhaomin:5:4}
  {\scriptsize}
  {\mbox{}}

\clearpage

\plotframe{assets/levels/zhaomin/level_05/frame_05.jpg}
  {赵敏线第五拍《弥勒佛庙》第5帧}
  {fig:level:zhaomin:5:5}
  {\scriptsize}
  {\mbox{}}

\subsubsection{第六拍·濠州夺婚}\label{plot:zhaomin:6}
本关以选择判断完成叙事向健康知识的转场，题目见\hyperref[question:zhaomin:6]{第\ref*{question:zhaomin:6}项}。

\plotframe{assets/levels/zhaomin/level_06/frame_01.jpg}
  {赵敏线第六拍《濠州夺婚》第1帧}
  {fig:level:zhaomin:6:1}
  {\scriptsize}
  {\textbf{旁白：}濠州城中，张无忌与周芷若大婚。喜堂之上，红绸高挂，宾客满座。范遥却见赵敏只身而来，他劝道：世上不如意事十居八九，既已如此，也是勉强不来了。\par}

\plotframe{assets/levels/zhaomin/level_06/frame_02.jpg}
  {赵敏线第六拍《濠州夺婚》第2帧}
  {fig:level:zhaomin:6:2}
  {\scriptsize}
  {\textbf{赵敏：}（微微一笑）我偏要勉强。\par
\textbf{旁白：}赵敏携金毛狮王谢逊的头发而来——张无忌的义父落在成昆手中，唯有她知道下落。这一声偏要勉强，成为金庸武侠中最动人的宣言。\par}

\clearpage

\plotframe{assets/levels/zhaomin/level_06/frame_03.jpg}
  {赵敏线第六拍《濠州夺婚》第3帧}
  {fig:level:zhaomin:6:3}
  {\scriptsize}
  {\textbf{周芷若：}（素手裂红裳，声音颤抖）张无忌！你若随她去了，我周芷若……此生……与你恩断义绝！\par}

\plotframe{assets/levels/zhaomin/level_06/frame_04.jpg}
  {赵敏线第六拍《濠州夺婚》第4帧}
  {fig:level:zhaomin:6:4}
  {\scriptsize}
  {\textbf{旁白：}婚礼变局，众人错愕。而张无忌最终还是随着赵敏走了。周芷若承受的打击，不仅是感情上的失落，更是精神上的重创。大喜大悲之间，最伤女子身心。\par}

\plotframe{assets/levels/zhaomin/level_06/frame_05.jpg}
  {赵敏线第六拍《濠州夺婚》第5帧}
  {fig:level:zhaomin:6:5}
  {\scriptsize}
  {\mbox{}}

\clearpage

\subsubsection{第七拍·少室山下}\label{plot:zhaomin:7}
本关以选择判断完成叙事向健康知识的转场，题目见\hyperref[question:zhaomin:7]{第\ref*{question:zhaomin:7}项}。

\plotframe{assets/levels/zhaomin/level_07/frame_01.jpg}
  {赵敏线第七拍《少室山下》第1帧}
  {fig:level:zhaomin:7:1}
  {\scriptsize}
  {\textbf{旁白：}少林屠狮大会在即，天下英雄齐聚少室山。赵敏随张无忌一路奔波，体力消耗巨大。这日，少林斋堂之中……\par}

\plotframe{assets/levels/zhaomin/level_07/frame_02.jpg}
  {赵敏线第七拍《少室山下》第2帧}
  {fig:level:zhaomin:7:2}
  {\scriptsize}
  {\textbf{张无忌：}我在蝴蝶谷学医时，胡青牛曾言：女子以血为本，补血先须健脾胃。脾胃为气血生化之源，若脾胃虚弱，纵食山珍海味亦难化生气血。\par}

\plotframe{assets/levels/zhaomin/level_07/frame_03.jpg}
  {赵敏线第七拍《少室山下》第3帧}
  {fig:level:zhaomin:7:3}
  {\scriptsize}
  {\textbf{赵敏：}无忌，连日赶路，我觉着浑身乏力。这些时日饮食不周，恐是气血两虚了。你可有什么食补之法？\par}

\clearpage

\plotframe{assets/levels/zhaomin/level_07/frame_04.jpg}
  {赵敏线第七拍《少室山下》第4帧}
  {fig:level:zhaomin:7:4}
  {\scriptsize}
  {\mbox{}}

\plotframe{assets/levels/zhaomin/level_07/frame_05.jpg}
  {赵敏线第七拍《少室山下》第5帧}
  {fig:level:zhaomin:7:5}
  {\scriptsize}
  {\mbox{}}

\subsubsection{第八拍·朔漠风霜}\label{plot:zhaomin:8}
本关以选择判断完成叙事向健康知识的转场，题目见\hyperref[question:zhaomin:8]{第\ref*{question:zhaomin:8}项}。

\plotframe{assets/levels/zhaomin/level_08/frame_01.jpg}
  {赵敏线第八拍《朔漠风霜》第1帧}
  {fig:level:zhaomin:8:1}
  {\scriptsize}
  {\textbf{旁白：}赵敏带张无忌回到蒙古草原。朔风如刀，大漠烈日照晒。赵敏虽是蒙古人，但多年在江南生活后重返朔漠，肌肤也需重新适应。\par}

\clearpage

\plotframe{assets/levels/zhaomin/level_08/frame_02.jpg}
  {赵敏线第八拍《朔漠风霜》第2帧}
  {fig:level:zhaomin:8:2}
  {\scriptsize}
  {\textbf{张无忌：}敏敏，这朔漠的风沙可比江南的烟雨厉害得多。你脸上都吹红了。\par}

\plotframe{assets/levels/zhaomin/level_08/frame_03.jpg}
  {赵敏线第八拍《朔漠风霜》第3帧}
  {fig:level:zhaomin:8:3}
  {\scriptsize}
  {\textbf{赵敏：}这有什么。我蒙古女子自有护肤之法——羊脂混沙棘油，再加些红花，涂于面上，可防风防晒。倒是你们中原女子，整日以轻纱遮面，以铅粉敷脸——我听说那铅粉有毒，可是真的？\par}

\plotframe{assets/levels/zhaomin/level_08/frame_04.jpg}
  {赵敏线第八拍《朔漠风霜》第4帧}
  {fig:level:zhaomin:8:4}
  {\scriptsize}
  {\mbox{}}

\clearpage

\plotframe{assets/levels/zhaomin/level_08/frame_05.jpg}
  {赵敏线第八拍《朔漠风霜》第5帧}
  {fig:level:zhaomin:8:5}
  {\scriptsize}
  {\mbox{}}

\subsubsection{第九拍·冰火之岛}\label{plot:zhaomin:9}
本关以多轮简答完成叙事向健康知识的转场，题目见\hyperref[question:zhaomin:9]{第\ref*{question:zhaomin:9}项}。

\plotframe{assets/levels/zhaomin/level_09/frame_01.jpg}
  {赵敏线第九拍《冰火之岛》第1帧}
  {fig:level:zhaomin:9:1}
  {\scriptsize}
  {\textbf{旁白：}赵敏与张无忌决定寻找他出生的地方——冰火岛。这极北之地，寒冷与炎热并存，对人的身体是极大考验。\par}

\plotframe{assets/levels/zhaomin/level_09/frame_02.jpg}
  {赵敏线第九拍《冰火之岛》第2帧}
  {fig:level:zhaomin:9:2}
  {\scriptsize}
  {\textbf{张无忌：}《黄帝内经》有云：人与天地相参也，与日月相应也。环境之寒热燥湿，皆可影响人体阴阳平衡。女子体质本偏阴柔，更需注意适应环境变迁。\par}

\clearpage

\plotframe{assets/levels/zhaomin/level_09/frame_03.jpg}
  {赵敏线第九拍《冰火之岛》第3帧}
  {fig:level:zhaomin:9:3}
  {\scriptsize}
  {\textbf{赵敏：}无忌，此处气候如此极端——冰山与温泉共存。人在此等环境中生活，体内阴阳必然大乱。你可知中医如何理解环境致病？\par}

\plotframe{assets/levels/zhaomin/level_09/frame_04.jpg}
  {赵敏线第九拍《冰火之岛》第4帧}
  {fig:level:zhaomin:9:4}
  {\scriptsize}
  {\mbox{}}

\plotframe{assets/levels/zhaomin/level_09/frame_05.jpg}
  {赵敏线第九拍《冰火之岛》第5帧}
  {fig:level:zhaomin:9:5}
  {\scriptsize}
  {\mbox{}}

\clearpage

\subsubsection{第十拍·父女恩仇}\label{plot:zhaomin:10}
本关以选择判断完成叙事向健康知识的转场，题目见\hyperref[question:zhaomin:10]{第\ref*{question:zhaomin:10}项}。

\plotframe{assets/levels/zhaomin/level_10/frame_01.jpg}
  {赵敏线第十拍《父女恩仇》第1帧}
  {fig:level:zhaomin:10:1}
  {\tiny}
  {\textbf{旁白：}赵敏带张无忌回蒙古见父王。她本以为血浓于水，却不想父兄早已知晓她与明教中人来往，雷霆大怒。\par
\textbf{汝阳王：}敏敏！你可知罪？身为郡主，竟与叛逆为伍！你若执意跟他走——便是与我大元为敌！与我汝阳王府恩断义绝！\par
\textbf{旁白：}赵敏此言，触及了中老年健康的重要议题。女性会经历围绝经期；男性若出现迟发性性腺功能减退，则需结合症状与实验室检查诊断，不能笼统称为人人都会经历的男性更年期。中年男女也都需要重视睡眠、运动、慢病筛查与压力管理。\par}

\plotframe{assets/levels/zhaomin/level_10/frame_02.jpg}
  {赵敏线第十拍《父女恩仇》第2帧}
  {fig:level:zhaomin:10:2}
  {\scriptsize}
  {\textbf{赵敏：}父王……女儿不孝。但女儿心意已决。只是——父王，您鬓边的白发又多了一些。女儿虽走，仍挂念父王身体。父王年事渐高，肝火易旺，少饮酒，少动怒才是……\par}

\plotframe{assets/levels/zhaomin/level_10/frame_03.jpg}
  {赵敏线第十拍《父女恩仇》第3帧}
  {fig:level:zhaomin:10:3}
  {\scriptsize}
  {\mbox{}}

\clearpage

\plotframe{assets/levels/zhaomin/level_10/frame_04.jpg}
  {赵敏线第十拍《父女恩仇》第4帧}
  {fig:level:zhaomin:10:4}
  {\scriptsize}
  {\mbox{}}

\plotframe{assets/levels/zhaomin/level_10/frame_05.jpg}
  {赵敏线第十拍《父女恩仇》第5帧}
  {fig:level:zhaomin:10:5}
  {\scriptsize}
  {\mbox{}}

\subsubsection{第十一拍·玄冥寒掌}\label{plot:zhaomin:11}
本关以选择判断完成叙事向健康知识的转场，题目见\hyperref[question:zhaomin:11]{第\ref*{question:zhaomin:11}项}。

\plotframe{assets/levels/zhaomin/level_11/frame_01.jpg}
  {赵敏线第十一拍《玄冥寒掌》第1帧}
  {fig:level:zhaomin:11:1}
  {\scriptsize}
  {\mbox{}}

\clearpage

\plotframe{assets/levels/zhaomin/level_11/frame_02.jpg}
  {赵敏线第十一拍《玄冥寒掌》第2帧}
  {fig:level:zhaomin:11:2}
  {\scriptsize}
  {\textbf{旁白：}张无忌幼年身中玄冥神掌，寒毒入体，命悬一线。胡青牛以毕生所学为他续命——这寒热辨证，正是中医八纲辨证的核心内容。\par}

\plotframe{assets/levels/zhaomin/level_11/frame_03.jpg}
  {赵敏线第十一拍《玄冥寒掌》第3帧}
  {fig:level:zhaomin:11:3}
  {\scriptsize}
  {\mbox{}}

\plotframe{assets/levels/zhaomin/level_11/frame_04.jpg}
  {赵敏线第十一拍《玄冥寒掌》第4帧}
  {fig:level:zhaomin:11:4}
  {\scriptsize}
  {\mbox{}}

\clearpage

\plotframe{assets/levels/zhaomin/level_11/frame_05.jpg}
  {赵敏线第十一拍《玄冥寒掌》第5帧}
  {fig:level:zhaomin:11:5}
  {\scriptsize}
  {\textbf{赵敏：}无忌曾说——玄冥寒掌之毒，发作时如坠冰窖，四肢百骸无不寒冷彻骨。我虽不懂武功，但知中医有寒者热之之法。他用九阳神功御寒，是不是便是此理？\par}

\subsubsection{第十二拍·大都宫闱}\label{plot:zhaomin:12}
本关以选择判断完成叙事向健康知识的转场，题目见\hyperref[question:zhaomin:12]{第\ref*{question:zhaomin:12}项}。

\plotframe{assets/levels/zhaomin/level_12/frame_01.jpg}
  {赵敏线第十二拍《大都宫闱》第1帧}
  {fig:level:zhaomin:12:1}
  {\scriptsize}
  {\textbf{旁白：}大都皇宫中，赵敏自幼耳濡目染——蒙古贵族争权夺利，毒药暗杀时有发生。她虽不愿参与，却也学会了辨识各种毒物。\par}

\plotframe{assets/levels/zhaomin/level_12/frame_02.jpg}
  {赵敏线第十二拍《大都宫闱》第2帧}
  {fig:level:zhaomin:12:2}
  {\scriptsize}
  {\textbf{御医：}郡主问得好。乌头大热，有剧毒，然炮制得法、配伍得当，却是治疗风寒湿痹之要药。正如《神农本草经》所言——药有酸咸甘苦辛五味，又有寒热温凉四气，及有毒无毒。世间万物，用之得当即为良药，用之失当即为毒药。\par}

\clearpage

\plotframe{assets/levels/zhaomin/level_12/frame_03.jpg}
  {赵敏线第十二拍《大都宫闱》第3帧}
  {fig:level:zhaomin:12:3}
  {\scriptsize}
  {\mbox{}}

\plotframe{assets/levels/zhaomin/level_12/frame_04.jpg}
  {赵敏线第十二拍《大都宫闱》第4帧}
  {fig:level:zhaomin:12:4}
  {\scriptsize}
  {\textbf{赵敏：}（对医师）先生，我父王府中有味药唤作乌头——据说治风湿有奇效，然用量稍过便能致命。此药究竟有毒无毒？\par}

\plotframe{assets/levels/zhaomin/level_12/frame_05.jpg}
  {赵敏线第十二拍《大都宫闱》第5帧}
  {fig:level:zhaomin:12:5}
  {\scriptsize}
  {\mbox{}}

\clearpage

\subsubsection{第十三拍·蝴蝶谷中}\label{plot:zhaomin:13}
本关以多轮简答完成叙事向健康知识的转场，题目见\hyperref[question:zhaomin:13]{第\ref*{question:zhaomin:13}项}。

\plotframe{assets/levels/zhaomin/level_13/frame_01.jpg}
  {赵敏线第十三拍《蝴蝶谷中》第1帧}
  {fig:level:zhaomin:13:1}
  {\scriptsize}
  {\textbf{旁白：}蝴蝶谷——蝶谷医仙胡青牛隐居之所。当年他曾在此救治张无忌，今日赵敏踏访遗址，探寻医道真谛。\par}

\plotframe{assets/levels/zhaomin/level_13/frame_02.jpg}
  {赵敏线第十三拍《蝴蝶谷中》第2帧}
  {fig:level:zhaomin:13:2}
  {\scriptsize}
  {\textbf{赵敏：}（抚摸着褪色的医书）胡先生当年教无忌——望闻问切四诊合参。我虽不学医，却也好奇：这望字，究竟能看出什么门道？\par}

\plotframe{assets/levels/zhaomin/level_13/frame_03.jpg}
  {赵敏线第十三拍《蝴蝶谷中》第3帧}
  {fig:level:zhaomin:13:3}
  {\scriptsize}
  {\textbf{张无忌：}胡师教我，望诊可察面色与舌象，但须四诊合参。唇甲苍白、面色异常或舌象变化，只能算线索，不能凭一眼便断病；若伴头晕、气促、发热或持续不适，还须验血与求医。望而知之虽称神，妄下断语却非医道。\par}

\clearpage

\plotframe{assets/levels/zhaomin/level_13/frame_04.jpg}
  {赵敏线第十三拍《蝴蝶谷中》第4帧}
  {fig:level:zhaomin:13:4}
  {\scriptsize}
  {\mbox{}}

\plotframe{assets/levels/zhaomin/level_13/frame_05.jpg}
  {赵敏线第十三拍《蝴蝶谷中》第5帧}
  {fig:level:zhaomin:13:5}
  {\scriptsize}
  {\mbox{}}

\subsubsection{第十四拍·光明顶上}\label{plot:zhaomin:14}
本关以选择判断完成叙事向健康知识的转场，题目见\hyperref[question:zhaomin:14]{第\ref*{question:zhaomin:14}项}。

\plotframe{assets/levels/zhaomin/level_14/frame_01.jpg}
  {赵敏线第十四拍《光明顶上》第1帧}
  {fig:level:zhaomin:14:1}
  {\scriptsize}
  {\textbf{旁白：}光明顶大战——张无忌以一己之力独战六大门派。乾坤大挪移虽威力无穷，但对身体消耗极大。\par}

\clearpage

\plotframe{assets/levels/zhaomin/level_14/frame_02.jpg}
  {赵敏线第十四拍《光明顶上》第2帧}
  {fig:level:zhaomin:14:2}
  {\scriptsize}
  {\mbox{}}

\plotframe{assets/levels/zhaomin/level_14/frame_03.jpg}
  {赵敏线第十四拍《光明顶上》第3帧}
  {fig:level:zhaomin:14:3}
  {\scriptsize}
  {\mbox{}}

\plotframe{assets/levels/zhaomin/level_14/frame_04.jpg}
  {赵敏线第十四拍《光明顶上》第4帧}
  {fig:level:zhaomin:14:4}
  {\scriptsize}
  {\textbf{张无忌：}你倒会比喻。不过确有相似之处——都是体力极度消耗，气血大亏。练武之人，尤其女子习武，更需注意运动适度，防止损伤经络筋骨。\par}

\clearpage

\plotframe{assets/levels/zhaomin/level_14/frame_05.jpg}
  {赵敏线第十四拍《光明顶上》第5帧}
  {fig:level:zhaomin:14:5}
  {\scriptsize}
  {\textbf{赵敏：}无忌，你这乾坤大挪移虽厉害，但每用一次都大汗淋漓、精疲力尽——这岂非与女子产后大伤元气相似？\par}

\subsubsection{第十五拍·波斯异术}\label{plot:zhaomin:15}
本关以选择判断完成叙事向健康知识的转场，题目见\hyperref[question:zhaomin:15]{第\ref*{question:zhaomin:15}项}。

\plotframe{assets/levels/zhaomin/level_15/frame_01.jpg}
  {赵敏线第十五拍《波斯异术》第1帧}
  {fig:level:zhaomin:15:1}
  {\scriptsize}
  {\textbf{旁白：}小昭远赴波斯继任明教教主。赵敏与张无忌收到她托商队带来的书信和异域药材……\par}

\plotframe{assets/levels/zhaomin/level_15/frame_02.jpg}
  {赵敏线第十五拍《波斯异术》第2帧}
  {fig:level:zhaomin:15:2}
  {\scriptsize}
  {\textbf{赵敏：}（拆开信）小昭说波斯医术与中原不同——他们以四液（血液、黏液、黄胆汁、黑胆汁）解释人体，倒是与我蒙古萨满的三体液之说有几分契合……还有一味唤作番红花的药材，说对女子经期颇有裨益。\par}

\clearpage

\plotframe{assets/levels/zhaomin/level_15/frame_03.jpg}
  {赵敏线第十五拍《波斯异术》第3帧}
  {fig:level:zhaomin:15:3}
  {\scriptsize}
  {\mbox{}}

\plotframe{assets/levels/zhaomin/level_15/frame_04.jpg}
  {赵敏线第十五拍《波斯异术》第4帧}
  {fig:level:zhaomin:15:4}
  {\scriptsize}
  {\mbox{}}

\plotframe{assets/levels/zhaomin/level_15/frame_05.jpg}
  {赵敏线第十五拍《波斯异术》第5帧}
  {fig:level:zhaomin:15:5}
  {\scriptsize}
  {\mbox{}}

\clearpage

\subsubsection{第十六拍·襄阳遗梦}\label{plot:zhaomin:16}
本关以选择判断完成叙事向健康知识的转场，题目见\hyperref[question:zhaomin:16]{第\ref*{question:zhaomin:16}项}。

\plotframe{assets/levels/zhaomin/level_16/frame_01.jpg}
  {赵敏线第十六拍《襄阳遗梦》第1帧}
  {fig:level:zhaomin:16:1}
  {\scriptsize}
  {\textbf{旁白：}途经襄阳——这座曾由大侠郭靖、黄蓉镇守数十年的城池，见证了宋元交替的血雨腥风。赵敏身为蒙古郡主，面对这座被蒙古大军攻陷的城池，心绪复杂。\par}

\plotframe{assets/levels/zhaomin/level_16/frame_02.jpg}
  {赵敏线第十六拍《襄阳遗梦》第2帧}
  {fig:level:zhaomin:16:2}
  {\scriptsize}
  {\mbox{}}

\plotframe{assets/levels/zhaomin/level_16/frame_03.jpg}
  {赵敏线第十六拍《襄阳遗梦》第3帧}
  {fig:level:zhaomin:16:3}
  {\scriptsize}
  {\mbox{}}

\clearpage

\plotframe{assets/levels/zhaomin/level_16/frame_04.jpg}
  {赵敏线第十六拍《襄阳遗梦》第4帧}
  {fig:level:zhaomin:16:4}
  {\scriptsize}
  {\textbf{赵敏：}（抚着城砖，叹息道）郭大侠、黄女侠以死守城……乱世之中，女子在战火中最是受难。无忌，若有一日天下大乱，你教我的急救之法——止血、包扎、正骨——或许能救下几条人命。\par}

\plotframe{assets/levels/zhaomin/level_16/frame_05.jpg}
  {赵敏线第十六拍《襄阳遗梦》第5帧}
  {fig:level:zhaomin:16:5}
  {\scriptsize}
  {\mbox{}}

\subsubsection{第十七拍·终南问道}\label{plot:zhaomin:17}
本关以多轮简答完成叙事向健康知识的转场，题目见\hyperref[question:zhaomin:17]{第\ref*{question:zhaomin:17}项}。

\plotframe{assets/levels/zhaomin/level_17/frame_01.jpg}
  {赵敏线第十七拍《终南问道》第1帧}
  {fig:level:zhaomin:17:1}
  {\scriptsize}
  {\textbf{旁白：}终南山——当年杨过与小龙女隐居之地。赵敏寻访古墓遗址，问道于山中隐士。\par}

\clearpage

\plotframe{assets/levels/zhaomin/level_17/frame_02.jpg}
  {赵敏线第十七拍《终南问道》第2帧}
  {fig:level:zhaomin:17:2}
  {\scriptsize}
  {\mbox{}}

\plotframe{assets/levels/zhaomin/level_17/frame_03.jpg}
  {赵敏线第十七拍《终南问道》第3帧}
  {fig:level:zhaomin:17:3}
  {\scriptsize}
  {\textbf{隐士：}郡主好问。古人谈养阴血、疏肝气、固肾精，落到日用，不过是饮食有节、起居有常、劳逸相济、情志有托。药膳与针灸皆有适应证和禁忌，不可见方就补、见穴就灸；若月事异常、久痛不愈，仍当求医。\par}

\plotframe{assets/levels/zhaomin/level_17/frame_04.jpg}
  {赵敏线第十七拍《终南问道》第4帧}
  {fig:level:zhaomin:17:4}
  {\scriptsize}
  {\textbf{赵敏：}道长远在终南修行，可知女子养生最重什么？我读了些医书，说是女子以血为本，以肝为先天——可这与养生有何关联？\par}

\clearpage

\plotframe{assets/levels/zhaomin/level_17/frame_05.jpg}
  {赵敏线第十七拍《终南问道》第5帧}
  {fig:level:zhaomin:17:5}
  {\scriptsize}
  {\mbox{}}

\subsubsection{第十八拍·天下归心}\label{plot:zhaomin:18}
本关以选择判断完成叙事向健康知识的转场，题目见\hyperref[question:zhaomin:18]{第\ref*{question:zhaomin:18}项}。

\plotframe{assets/levels/zhaomin/level_18/frame_01.jpg}
  {赵敏线第十八拍《天下归心》第1帧}
  {fig:level:zhaomin:18:1}
  {\scriptsize}
  {\textbf{旁白：}十八道关卡，十八次考验。赵敏从蒙古郡主，到明教教主夫人，再到冰火岛的归隐者——她的每一步，都走在自己选择的路上。\par
\textbf{赵敏：}（微微一笑）你倒提醒了我。我虽生在王侯之家，却见过太多民间女子的苦楚——经期无人关照、产后无人照料、老来孤苦伶仃……若有此书，愿她们能多爱惜自己一些。\par}

\plotframe{assets/levels/zhaomin/level_18/frame_02.jpg}
  {赵敏线第十八拍《天下归心》第2帧}
  {fig:level:zhaomin:18:2}
  {\scriptsize}
  {\textbf{张无忌：}敏敏，你从不会医术，到如今通晓女子养生之道——若有一日，能将这十八关中所学的医道整理成书，传给天下女子，何尝不是一件功德？\par}

\clearpage

\plotframe{assets/levels/zhaomin/level_18/frame_03.jpg}
  {赵敏线第十八拍《天下归心》第3帧}
  {fig:level:zhaomin:18:3}
  {\scriptsize}
  {\mbox{}}

\plotframe{assets/levels/zhaomin/level_18/frame_04.jpg}
  {赵敏线第十八拍《天下归心》第4帧}
  {fig:level:zhaomin:18:4}
  {\scriptsize}
  {\textbf{赵敏：}（望着大漠夕阳，轻声道）无忌，我这一生，最庆幸的——便是当年在绿柳山庄地牢中遇见了你。那些江湖恩怨、家国纷争，都已成了前尘往事。如今只愿——天下太平，人间无病。\par}

\plotframe{assets/levels/zhaomin/level_18/frame_05.jpg}
  {赵敏线第十八拍《天下归心》第5帧}
  {fig:level:zhaomin:18:5}
  {\scriptsize}
  {\mbox{}}

\clearpage
\subsubsection{路线结局：第十八拍 · 天下归心}
\textbf{结尾诗}\begin{quote}朔漠风烟散，天涯共月明。\par 胡笳声渐远，人间无病声。\par 冰火同归处，牛羊自在行。\par 愿得四海平，不闻战鼓鸣。\end{quote}
\textbf{结局文本}\begin{quote}赵敏与张无忌归隐冰火岛。她曾是大元郡主，执掌天下兵马；如今只愿在冰与火的交界处，牧羊养马，与心爱之人共度余生。\par \par 毡房外，大漠夕阳如血。她常想起父兄——愿蒙古草原风调雨顺，愿中原大地国泰民安，愿天下女子皆得健康，愿世间再无战乱与病痛。\par \par 蔡文姬作《胡笳十八拍》，叹乱世流离、母子永诀。千年之后，赵敏已不需经历那般悲苦——但那份对家国天下的关怀，对人间太平的祈愿，古今一同。\end{quote}

\clearpage
\subsection{周芷若线：汉水白莲}

\subsubsection{第一拍·汉水初遇}\label{plot:zhouzhiruo:1}
本关以选择判断完成叙事向健康知识的转场，题目见\hyperref[question:zhouzhiruo:1]{第\ref*{question:zhouzhiruo:1}项}。

\plotframe{assets/levels/zhouzhiruo/level_01/frame_01.jpg}
  {周芷若线第一拍《汉水初遇》第1帧}
  {fig:level:zhouzhiruo:1:1}
  {\scriptsize}
  {\textbf{旁白：}元末。汉水之上，一叶扁舟。渔家女周芷若的父亲病重垂危，张三丰携张无忌路过，见者心酸。\par
\textbf{旁白：}周芷若的父亲死于长期营养不良与感染。在元末战乱年代，百姓流离失所，饥饿是最常见的疾病。对女性而言，营养不良导致的贫血尤为普遍——面色苍白、头晕乏力、月经量少甚至闭经。\par}

\plotframe{assets/levels/zhouzhiruo/level_01/frame_02.jpg}
  {周芷若线第一拍《汉水初遇》第2帧}
  {fig:level:zhouzhiruo:1:2}
  {\scriptsize}
  {\textbf{周芷若：}（跪在父亲身旁，泪如雨下）爹爹……你不要走……芷若只有你了……\par}

\plotframe{assets/levels/zhouzhiruo/level_01/frame_03.jpg}
  {周芷若线第一拍《汉水初遇》第3帧}
  {fig:level:zhouzhiruo:1:3}
  {\scriptsize}
  {\mbox{}}

\clearpage

\plotframe{assets/levels/zhouzhiruo/level_01/frame_04.jpg}
  {周芷若线第一拍《汉水初遇》第4帧}
  {fig:level:zhouzhiruo:1:4}
  {\scriptsize}
  {\textbf{张三丰：}（诊脉后，叹息一声）小姑娘，你父亲是长期饥饿劳苦、气血两亏，又感染风寒，肺气已绝——贫道也无能为力了。\par}

\plotframe{assets/levels/zhouzhiruo/level_01/frame_05.jpg}
  {周芷若线第一拍《汉水初遇》第5帧}
  {fig:level:zhouzhiruo:1:5}
  {\scriptsize}
  {\mbox{}}

\subsubsection{第二拍·峨眉入门}\label{plot:zhouzhiruo:2}
本关以选择判断完成叙事向健康知识的转场，题目见\hyperref[question:zhouzhiruo:2]{第\ref*{question:zhouzhiruo:2}项}。

\plotframe{assets/levels/zhouzhiruo/level_02/frame_01.jpg}
  {周芷若线第二拍《峨眉入门》第1帧}
  {fig:level:zhouzhiruo:2:1}
  {\scriptsize}
  {\textbf{旁白：}周芷若被张三丰送至峨眉山，拜入灭绝师太门下。从垂髫少女到亭亭玉立，她在峨眉山度过了整个青春期。\par
\textbf{周芷若：}（恭敬地）徒儿不知，请师父明示。\par}

\clearpage

\plotframe{assets/levels/zhouzhiruo/level_02/frame_02.jpg}
  {周芷若线第二拍《峨眉入门》第2帧}
  {fig:level:zhouzhiruo:2:2}
  {\scriptsize}
  {\mbox{}}

\plotframe{assets/levels/zhouzhiruo/level_02/frame_03.jpg}
  {周芷若线第二拍《峨眉入门》第3帧}
  {fig:level:zhouzhiruo:2:3}
  {\tiny}
  {\textbf{灭绝师太：}芷若，你既入我峨眉门下，当知女子习武与男子不同。女子有经、带、胎、产之生理，练功需顺应天时，不可妄自逞强。你可知何为天癸？\par
\textbf{灭绝师太：}《黄帝内经》曰：女子二七而天癸至，任脉通，太冲脉盛，月事以时下。天癸者，乃肾中精气所化，主生殖发育。你来峨眉时方才十岁，如今十四，天癸将至——此后炼功，经期须减量，忌寒凉、忌过劳。此为女子养身之本，万不可轻忽！\par}

\plotframe{assets/levels/zhouzhiruo/level_02/frame_04.jpg}
  {周芷若线第二拍《峨眉入门》第4帧}
  {fig:level:zhouzhiruo:2:4}
  {\scriptsize}
  {\mbox{}}

\clearpage

\plotframe{assets/levels/zhouzhiruo/level_02/frame_05.jpg}
  {周芷若线第二拍《峨眉入门》第5帧}
  {fig:level:zhouzhiruo:2:5}
  {\scriptsize}
  {\mbox{}}

\subsubsection{第三拍·光明一剑}\label{plot:zhouzhiruo:3}
本关以选择判断完成叙事向健康知识的转场，题目见\hyperref[question:zhouzhiruo:3]{第\ref*{question:zhouzhiruo:3}项}。

\plotframe{assets/levels/zhouzhiruo/level_03/frame_01.jpg}
  {周芷若线第三拍《光明一剑》第1帧}
  {fig:level:zhouzhiruo:3:1}
  {\scriptsize}
  {\textbf{旁白：}光明顶上，六大门派围攻明教。灭绝师太命周芷若用倚天剑刺杀张无忌——那个曾在汉水之上对她有一饭之恩的少年。\par}

\plotframe{assets/levels/zhouzhiruo/level_03/frame_02.jpg}
  {周芷若线第三拍《光明一剑》第2帧}
  {fig:level:zhouzhiruo:3:2}
  {\scriptsize}
  {\textbf{周芷若：}（手颤抖，泪流满面）师父……我……\par}

\clearpage

\plotframe{assets/levels/zhouzhiruo/level_03/frame_03.jpg}
  {周芷若线第三拍《光明一剑》第3帧}
  {fig:level:zhouzhiruo:3:3}
  {\scriptsize}
  {\textbf{灭绝师太：}芷若！刺下去！此人是魔教教主，人人得而诛之！你若违抗师命，便不是我峨眉弟子！\par}

\plotframe{assets/levels/zhouzhiruo/level_03/frame_04.jpg}
  {周芷若线第三拍《光明一剑》第4帧}
  {fig:level:zhouzhiruo:3:4}
  {\scriptsize}
  {\textbf{旁白：}周芷若一剑刺出——刺入了张无忌的胸膛。这一剑，也在她自己心上留下了永久的伤口。在极度压力下做出的行为，往往会在日后造成创伤后应激——那些在巨大压力下被逼迫做出违背本心之事的女性，尤需关注心理健康。\par}

\plotframe{assets/levels/zhouzhiruo/level_03/frame_05.jpg}
  {周芷若线第三拍《光明一剑》第5帧}
  {fig:level:zhouzhiruo:3:5}
  {\scriptsize}
  {\mbox{}}

\clearpage

\subsubsection{第四拍·万安寺誓}\label{plot:zhouzhiruo:4}
本关以多轮简答完成叙事向健康知识的转场，题目见\hyperref[question:zhouzhiruo:4]{第\ref*{question:zhouzhiruo:4}项}。

\plotframe{assets/levels/zhouzhiruo/level_04/frame_01.jpg}
  {周芷若线第四拍《万安寺誓》第1帧}
  {fig:level:zhouzhiruo:4:1}
  {\scriptsize}
  {\textbf{旁白：}万安寺大火。灭绝师太在临终前将周芷若叫到身边，逼她立下毒誓——\par}

\plotframe{assets/levels/zhouzhiruo/level_04/frame_02.jpg}
  {周芷若线第四拍《万安寺誓》第2帧}
  {fig:level:zhouzhiruo:4:2}
  {\scriptsize}
  {\textbf{灭绝师太：}芷若，你跪下！对为师发誓——第一，夺回倚天剑与屠龙刀，光复我峨眉派！第二，不得嫁给张无忌那个魔教妖人！若有违誓……\par}

\plotframe{assets/levels/zhouzhiruo/level_04/frame_03.jpg}
  {周芷若线第四拍《万安寺誓》第3帧}
  {fig:level:zhouzhiruo:4:3}
  {\scriptsize}
  {\textbf{周芷若：}（浑身发抖，声音嘶哑）徒儿……发誓……若有违誓……天诛地灭，不得好死……\par}

\clearpage

\plotframe{assets/levels/zhouzhiruo/level_04/frame_04.jpg}
  {周芷若线第四拍《万安寺誓》第4帧}
  {fig:level:zhouzhiruo:4:4}
  {\scriptsize}
  {\textbf{旁白：}在巨大精神压力下被迫立下的毒誓，如同梦魇一般缠绕周芷若。长期的精神压力可导致月经不调甚至闭经——中医称气滞血瘀，冲任不通。现代医学证实，下丘脑性闭经（FHA）与精神压力密切相关。\par}

\plotframe{assets/levels/zhouzhiruo/level_04/frame_05.jpg}
  {周芷若线第四拍《万安寺誓》第5帧}
  {fig:level:zhouzhiruo:4:5}
  {\scriptsize}
  {\mbox{}}

\subsubsection{第五拍·灵蛇毒计}\label{plot:zhouzhiruo:5}
本关以选择判断完成叙事向健康知识的转场，题目见\hyperref[question:zhouzhiruo:5]{第\ref*{question:zhouzhiruo:5}项}。

\plotframe{assets/levels/zhouzhiruo/level_05/frame_01.jpg}
  {周芷若线第五拍《灵蛇毒计》第1帧}
  {fig:level:zhouzhiruo:5:1}
  {\scriptsize}
  {\textbf{旁白：}灵蛇岛上，倚天剑与屠龙刀首次聚齐。周芷若趁众人熟睡之际，以十香软筋散迷倒众人，划伤殷离的面容，嫁祸赵敏……\par}

\clearpage

\plotframe{assets/levels/zhouzhiruo/level_05/frame_02.jpg}
  {周芷若线第五拍《灵蛇毒计》第2帧}
  {fig:level:zhouzhiruo:5:2}
  {\scriptsize}
  {\textbf{旁白：}十香软筋散——虽为小说中的虚构毒药，但软筋之意与现实中能导致肌肉松弛、呼吸抑制的神经毒素有相似之处。古今中外，毒物辨识始终是医学的重要分支。\par}

\plotframe{assets/levels/zhouzhiruo/level_05/frame_03.jpg}
  {周芷若线第五拍《灵蛇毒计》第3帧}
  {fig:level:zhouzhiruo:5:3}
  {\scriptsize}
  {\mbox{}}

\plotframe{assets/levels/zhouzhiruo/level_05/frame_04.jpg}
  {周芷若线第五拍《灵蛇毒计》第4帧}
  {fig:level:zhouzhiruo:5:4}
  {\scriptsize}
  {\mbox{}}

\clearpage

\plotframe{assets/levels/zhouzhiruo/level_05/frame_05.jpg}
  {周芷若线第五拍《灵蛇毒计》第5帧}
  {fig:level:zhouzhiruo:5:5}
  {\scriptsize}
  {\mbox{}}

\subsubsection{第六拍·裂裳之痛}\label{plot:zhouzhiruo:6}
本关以选择判断完成叙事向健康知识的转场，题目见\hyperref[question:zhouzhiruo:6]{第\ref*{question:zhouzhiruo:6}项}。

\plotframe{assets/levels/zhouzhiruo/level_06/frame_01.jpg}
  {周芷若线第六拍《裂裳之痛》第1帧}
  {fig:level:zhouzhiruo:6:1}
  {\scriptsize}
  {\textbf{旁白：}濠州城。张无忌与周芷若的婚礼。满堂宾客，红绸高挂。就在第三拜即将完成之时——赵敏出现了。\par
\textbf{张无忌：}（神色大变）义父……\par}

\plotframe{assets/levels/zhouzhiruo/level_06/frame_02.jpg}
  {周芷若线第六拍《裂裳之痛》第2帧}
  {fig:level:zhouzhiruo:6:2}
  {\scriptsize}
  {\textbf{赵敏：}张无忌！你义父谢逊的下落——只有我知道。你若今日拜了堂，今生休想再见他一面！\par}

\clearpage

\plotframe{assets/levels/zhouzhiruo/level_06/frame_03.jpg}
  {周芷若线第六拍《裂裳之痛》第3帧}
  {fig:level:zhouzhiruo:6:3}
  {\scriptsize}
  {\textbf{周芷若：}（凄然一笑，双手抓住嫁衣前襟，用力一撕——）既然如此——这婚，不成也罢！\par}

\plotframe{assets/levels/zhouzhiruo/level_06/frame_04.jpg}
  {周芷若线第六拍《裂裳之痛》第4帧}
  {fig:level:zhouzhiruo:6:4}
  {\scriptsize}
  {\textbf{旁白：}新妇素手裂红裳——这一撕，撕裂的不只是嫁衣，还有周芷若对爱情的最后一丝幻想。大喜变大悲，对女性身心的打击，不亚于一场重病。\par}

\plotframe{assets/levels/zhouzhiruo/level_06/frame_05.jpg}
  {周芷若线第六拍《裂裳之痛》第5帧}
  {fig:level:zhouzhiruo:6:5}
  {\scriptsize}
  {\mbox{}}

\clearpage

\subsubsection{第七拍·九阴白骨}\label{plot:zhouzhiruo:7}
本关以选择判断完成叙事向健康知识的转场，题目见\hyperref[question:zhouzhiruo:7]{第\ref*{question:zhouzhiruo:7}项}。

\plotframe{assets/levels/zhouzhiruo/level_07/frame_01.jpg}
  {周芷若线第七拍《九阴白骨》第1帧}
  {fig:level:zhouzhiruo:7:1}
  {\scriptsize}
  {\textbf{旁白：}失去爱情后的周芷若，将全部精力投入九阴白骨爪的修炼。但她不知道，过度训练正在摧毁她的双手——正如执念正在摧毁她的内心。\par}

\plotframe{assets/levels/zhouzhiruo/level_07/frame_02.jpg}
  {周芷若线第七拍《九阴白骨》第2帧}
  {fig:level:zhouzhiruo:7:2}
  {\scriptsize}
  {\mbox{}}

\plotframe{assets/levels/zhouzhiruo/level_07/frame_03.jpg}
  {周芷若线第七拍《九阴白骨》第3帧}
  {fig:level:zhouzhiruo:7:3}
  {\scriptsize}
  {\textbf{周芷若：}（揉着剧痛的指关节）这九阴白骨爪……每练一日，手指便痛上三分。可若不练，如何胜过赵敏？如何在武林中立足？\par}

\clearpage

\plotframe{assets/levels/zhouzhiruo/level_07/frame_04.jpg}
  {周芷若线第七拍《九阴白骨》第4帧}
  {fig:level:zhouzhiruo:7:4}
  {\scriptsize}
  {\mbox{}}

\plotframe{assets/levels/zhouzhiruo/level_07/frame_05.jpg}
  {周芷若线第七拍《九阴白骨》第5帧}
  {fig:level:zhouzhiruo:7:5}
  {\scriptsize}
  {\mbox{}}

\subsubsection{第八拍·峨眉掌门}\label{plot:zhouzhiruo:8}
本关以选择判断完成叙事向健康知识的转场，题目见\hyperref[question:zhouzhiruo:8]{第\ref*{question:zhouzhiruo:8}项}。

\plotframe{assets/levels/zhouzhiruo/level_08/frame_01.jpg}
  {周芷若线第八拍《峨眉掌门》第1帧}
  {fig:level:zhouzhiruo:8:1}
  {\scriptsize}
  {\textbf{周芷若：}（对镜自照，看到眼角的细纹）不过数月光景，却仿佛老了十年。掌门之职——比练九阴白骨爪更消磨人啊。\par}

\clearpage

\plotframe{assets/levels/zhouzhiruo/level_08/frame_02.jpg}
  {周芷若线第八拍《峨眉掌门》第2帧}
  {fig:level:zhouzhiruo:8:2}
  {\scriptsize}
  {\textbf{旁白：}周芷若成为峨眉掌门后，日夜操劳——既要处理派务，又要修炼武功，还要应对江湖风波。她常常半夜还在批阅文书，日渐憔悴。\par}

\plotframe{assets/levels/zhouzhiruo/level_08/frame_03.jpg}
  {周芷若线第八拍《峨眉掌门》第3帧}
  {fig:level:zhouzhiruo:8:3}
  {\scriptsize}
  {\mbox{}}

\plotframe{assets/levels/zhouzhiruo/level_08/frame_04.jpg}
  {周芷若线第八拍《峨眉掌门》第4帧}
  {fig:level:zhouzhiruo:8:4}
  {\scriptsize}
  {\mbox{}}

\clearpage

\plotframe{assets/levels/zhouzhiruo/level_08/frame_05.jpg}
  {周芷若线第八拍《峨眉掌门》第5帧}
  {fig:level:zhouzhiruo:8:5}
  {\scriptsize}
  {\mbox{}}

\subsubsection{第九拍·屠狮之会}\label{plot:zhouzhiruo:9}
本关以多轮简答完成叙事向健康知识的转场，题目见\hyperref[question:zhouzhiruo:9]{第\ref*{question:zhouzhiruo:9}项}。

\plotframe{assets/levels/zhouzhiruo/level_09/frame_01.jpg}
  {周芷若线第九拍《屠狮之会》第1帧}
  {fig:level:zhouzhiruo:9:1}
  {\scriptsize}
  {\textbf{旁白：}屠狮大会上，周芷若以九阴白骨爪横扫群雄。然而白天越是强大，夜晚越是恐惧——她已连续多日失眠，噩梦不断。\par}

\plotframe{assets/levels/zhouzhiruo/level_09/frame_02.jpg}
  {周芷若线第九拍《屠狮之会》第2帧}
  {fig:level:zhouzhiruo:9:2}
  {\scriptsize}
  {\mbox{}}

\clearpage

\plotframe{assets/levels/zhouzhiruo/level_09/frame_03.jpg}
  {周芷若线第九拍《屠狮之会》第3帧}
  {fig:level:zhouzhiruo:9:3}
  {\scriptsize}
  {\textbf{周芷若：}（深夜独白）每一个夜晚，我都会梦到师父……梦到灵蛇岛上殷离的脸……梦到濠州喜堂上那一幕……我何尝不想安睡？可一闭眼，便是这些……\par}

\plotframe{assets/levels/zhouzhiruo/level_09/frame_04.jpg}
  {周芷若线第九拍《屠狮之会》第4帧}
  {fig:level:zhouzhiruo:9:4}
  {\scriptsize}
  {\mbox{}}

\plotframe{assets/levels/zhouzhiruo/level_09/frame_05.jpg}
  {周芷若线第九拍《屠狮之会》第5帧}
  {fig:level:zhouzhiruo:9:5}
  {\scriptsize}
  {\mbox{}}

\clearpage

\subsubsection{第十拍·黄衫女子}\label{plot:zhouzhiruo:10}
本关以选择判断完成叙事向健康知识的转场，题目见\hyperref[question:zhouzhiruo:10]{第\ref*{question:zhouzhiruo:10}项}。

\plotframe{assets/levels/zhouzhiruo/level_10/frame_01.jpg}
  {周芷若线第十拍《黄衫女子》第1帧}
  {fig:level:zhouzhiruo:10:1}
  {\scriptsize}
  {\textbf{旁白：}黄衫女子出现——她以正宗九阴真经武功轻易击败周芷若的九阴白骨爪。她不仅是武功上的胜者，更是一种女性力量的象征——强大而慈悲，独立而从容。\par}

\plotframe{assets/levels/zhouzhiruo/level_10/frame_02.jpg}
  {周芷若线第十拍《黄衫女子》第2帧}
  {fig:level:zhouzhiruo:10:2}
  {\scriptsize}
  {\mbox{}}

\plotframe{assets/levels/zhouzhiruo/level_10/frame_03.jpg}
  {周芷若线第十拍《黄衫女子》第3帧}
  {fig:level:zhouzhiruo:10:3}
  {\scriptsize}
  {\textbf{黄衫女子：}周姑娘，你天赋极高，奈何心术不正。女子一生，不必依附于任何人——你执着于一个男子，可曾问过自己：没有他，你便不完整了吗？\par}

\clearpage

\plotframe{assets/levels/zhouzhiruo/level_10/frame_04.jpg}
  {周芷若线第十拍《黄衫女子》第4帧}
  {fig:level:zhouzhiruo:10:4}
  {\scriptsize}
  {\mbox{}}

\plotframe{assets/levels/zhouzhiruo/level_10/frame_05.jpg}
  {周芷若线第十拍《黄衫女子》第5帧}
  {fig:level:zhouzhiruo:10:5}
  {\scriptsize}
  {\mbox{}}

\subsubsection{第十一拍·问心有愧}\label{plot:zhouzhiruo:11}
本关以选择判断完成叙事向健康知识的转场，题目见\hyperref[question:zhouzhiruo:11]{第\ref*{question:zhouzhiruo:11}项}。

\plotframe{assets/levels/zhouzhiruo/level_11/frame_01.jpg}
  {周芷若线第十一拍《问心有愧》第1帧}
  {fig:level:zhouzhiruo:11:1}
  {\scriptsize}
  {\textbf{旁白：}当张无忌问周芷若是否还恨他时，她答了一句——这句话成为金庸武侠中最令人心碎的台词之一。\par}

\clearpage

\plotframe{assets/levels/zhouzhiruo/level_11/frame_02.jpg}
  {周芷若线第十一拍《问心有愧》第2帧}
  {fig:level:zhouzhiruo:11:2}
  {\scriptsize}
  {\textbf{周芷若：}（低声道）倘若我问心有愧呢？\par}

\plotframe{assets/levels/zhouzhiruo/level_11/frame_03.jpg}
  {周芷若线第十一拍《问心有愧》第3帧}
  {fig:level:zhouzhiruo:11:3}
  {\scriptsize}
  {\mbox{}}

\plotframe{assets/levels/zhouzhiruo/level_11/frame_04.jpg}
  {周芷若线第十一拍《问心有愧》第4帧}
  {fig:level:zhouzhiruo:11:4}
  {\scriptsize}
  {\textbf{旁白：}这句话里有多少执念、多少不甘、多少对自己的失望？《金匮要略》中有一方剂名甘麦大枣汤，专治妇人脏躁，悲伤欲哭，象如神灵所作——周芷若之症，与此何其相似。\par}

\clearpage

\plotframe{assets/levels/zhouzhiruo/level_11/frame_05.jpg}
  {周芷若线第十一拍《问心有愧》第5帧}
  {fig:level:zhouzhiruo:11:5}
  {\scriptsize}
  {\mbox{}}

\subsubsection{第十二拍·汉水归舟}\label{plot:zhouzhiruo:12}
本关以选择判断完成叙事向健康知识的转场，题目见\hyperref[question:zhouzhiruo:12]{第\ref*{question:zhouzhiruo:12}项}。

\plotframe{assets/levels/zhouzhiruo/level_12/frame_01.jpg}
  {周芷若线第十二拍《汉水归舟》第1帧}
  {fig:level:zhouzhiruo:12:1}
  {\scriptsize}
  {\textbf{旁白：}周芷若回到汉水——她的故乡。渔村中，有妇人正在坐月子，有新生儿在啼哭——这些最平常的生命场景，触动了她内心最柔软的地方。\par}

\plotframe{assets/levels/zhouzhiruo/level_12/frame_02.jpg}
  {周芷若线第十二拍《汉水归舟》第2帧}
  {fig:level:zhouzhiruo:12:2}
  {\scriptsize}
  {\textbf{周芷若：}（望着那产妇，轻声道）我母亲生我时，难产而死。我从未见过她……若当年有人懂些产科之术，或许——\par}

\clearpage

\plotframe{assets/levels/zhouzhiruo/level_12/frame_03.jpg}
  {周芷若线第十二拍《汉水归舟》第3帧}
  {fig:level:zhouzhiruo:12:3}
  {\scriptsize}
  {\mbox{}}

\plotframe{assets/levels/zhouzhiruo/level_12/frame_04.jpg}
  {周芷若线第十二拍《汉水归舟》第4帧}
  {fig:level:zhouzhiruo:12:4}
  {\scriptsize}
  {\mbox{}}

\plotframe{assets/levels/zhouzhiruo/level_12/frame_05.jpg}
  {周芷若线第十二拍《汉水归舟》第5帧}
  {fig:level:zhouzhiruo:12:5}
  {\scriptsize}
  {\mbox{}}

\clearpage

\subsubsection{第十三拍·峨眉金顶}\label{plot:zhouzhiruo:13}
本关以选择判断完成叙事向健康知识的转场，题目见\hyperref[question:zhouzhiruo:13]{第\ref*{question:zhouzhiruo:13}项}。

\plotframe{assets/levels/zhouzhiruo/level_13/frame_01.jpg}
  {周芷若线第十三拍《峨眉金顶》第1帧}
  {fig:level:zhouzhiruo:13:1}
  {\scriptsize}
  {\textbf{旁白：}峨眉金顶，周芷若遇到一位老居士——她弯腰驼背，每走一步都小心翼翼。骨质疏松，已悄悄夺走了她的身高和行动能力。\par}

\plotframe{assets/levels/zhouzhiruo/level_13/frame_02.jpg}
  {周芷若线第十三拍《峨眉金顶》第2帧}
  {fig:level:zhouzhiruo:13:2}
  {\scriptsize}
  {\textbf{老居士：}周掌门，老身年轻时也曾健步如飞。如今年过六十，背也弯了，膝盖也疼了，连抱个孙子都难……大夫说是骨质疏松。老身年轻时若有今日的见识，早些补钙、多晒太阳，也不至于此。\par}

\plotframe{assets/levels/zhouzhiruo/level_13/frame_03.jpg}
  {周芷若线第十三拍《峨眉金顶》第3帧}
  {fig:level:zhouzhiruo:13:3}
  {\scriptsize}
  {\mbox{}}

\clearpage

\plotframe{assets/levels/zhouzhiruo/level_13/frame_04.jpg}
  {周芷若线第十三拍《峨眉金顶》第4帧}
  {fig:level:zhouzhiruo:13:4}
  {\scriptsize}
  {\mbox{}}

\plotframe{assets/levels/zhouzhiruo/level_13/frame_05.jpg}
  {周芷若线第十三拍《峨眉金顶》第5帧}
  {fig:level:zhouzhiruo:13:5}
  {\scriptsize}
  {\mbox{}}

\subsubsection{第十四拍·药王谷中}\label{plot:zhouzhiruo:14}
本关以选择判断完成叙事向健康知识的转场，题目见\hyperref[question:zhouzhiruo:14]{第\ref*{question:zhouzhiruo:14}项}。

\plotframe{assets/levels/zhouzhiruo/level_14/frame_01.jpg}
  {周芷若线第十四拍《药王谷中》第1帧}
  {fig:level:zhouzhiruo:14:1}
  {\scriptsize}
  {\textbf{旁白：}周芷若行至传说中的药王谷——此处为唐代药王孙思邈晚年隐居采药之地（一说）。谷中药圃繁茂，本草典籍浩如烟海。\par}

\clearpage

\plotframe{assets/levels/zhouzhiruo/level_14/frame_02.jpg}
  {周芷若线第十四拍《药王谷中》第2帧}
  {fig:level:zhouzhiruo:14:2}
  {\scriptsize}
  {\textbf{药王谷老者：}周掌门既来求教，老夫便先从这四味药讲起——当归、白芍、熟地、川芎。此四味便是四物汤——妇科第一方。你可知道它们各自的作用？\par}

\plotframe{assets/levels/zhouzhiruo/level_14/frame_03.jpg}
  {周芷若线第十四拍《药王谷中》第3帧}
  {fig:level:zhouzhiruo:14:3}
  {\scriptsize}
  {\mbox{}}

\plotframe{assets/levels/zhouzhiruo/level_14/frame_04.jpg}
  {周芷若线第十四拍《药王谷中》第4帧}
  {fig:level:zhouzhiruo:14:4}
  {\scriptsize}
  {\mbox{}}

\clearpage

\plotframe{assets/levels/zhouzhiruo/level_14/frame_05.jpg}
  {周芷若线第十四拍《药王谷中》第5帧}
  {fig:level:zhouzhiruo:14:5}
  {\scriptsize}
  {\mbox{}}

\subsubsection{第十五拍·襄阳城墙}\label{plot:zhouzhiruo:15}
本关以多轮简答完成叙事向健康知识的转场，题目见\hyperref[question:zhouzhiruo:15]{第\ref*{question:zhouzhiruo:15}项}。

\plotframe{assets/levels/zhouzhiruo/level_15/frame_01.jpg}
  {周芷若线第十五拍《襄阳城墙》第1帧}
  {fig:level:zhouzhiruo:15:1}
  {\scriptsize}
  {\textbf{旁白：}襄阳城下，桃花盛开。周芷若在此偶遇一位正在求医的妇人——她乳房扪之有肿块，心中惶恐不安。\par}

\plotframe{assets/levels/zhouzhiruo/level_15/frame_02.jpg}
  {周芷若线第十五拍《襄阳城墙》第2帧}
  {fig:level:zhouzhiruo:15:2}
  {\scriptsize}
  {\mbox{}}

\clearpage

\plotframe{assets/levels/zhouzhiruo/level_15/frame_03.jpg}
  {周芷若线第十五拍《襄阳城墙》第3帧}
  {fig:level:zhouzhiruo:15:3}
  {\scriptsize}
  {\textbf{妇人：}周掌门……我这胸口有个硬块，不痛不痒的……我听说乳岩（乳腺癌）便是如此，心中日夜不安……\par
\textbf{周芷若：}（为妇人诊脉后，神色郑重）大姐莫慌。乳房肿块未必便是恶疾——有些是良性的乳癖（乳腺增生），有些需进一步检查方能断定。中医认为此与肝气郁结有关——你可知女子情志不畅，最易伤乳？\par}

\plotframe{assets/levels/zhouzhiruo/level_15/frame_04.jpg}
  {周芷若线第十五拍《襄阳城墙》第4帧}
  {fig:level:zhouzhiruo:15:4}
  {\scriptsize}
  {\mbox{}}

\plotframe{assets/levels/zhouzhiruo/level_15/frame_05.jpg}
  {周芷若线第十五拍《襄阳城墙》第5帧}
  {fig:level:zhouzhiruo:15:5}
  {\scriptsize}
  {\mbox{}}

\clearpage

\subsubsection{第十六拍·终南古墓}\label{plot:zhouzhiruo:16}
本关以选择判断完成叙事向健康知识的转场，题目见\hyperref[question:zhouzhiruo:16]{第\ref*{question:zhouzhiruo:16}项}。

\plotframe{assets/levels/zhouzhiruo/level_16/frame_01.jpg}
  {周芷若线第十六拍《终南古墓》第1帧}
  {fig:level:zhouzhiruo:16:1}
  {\scriptsize}
  {\textbf{旁白：}周芷若寻至终南山活死人墓——当年小龙女在古墓中独居数十年，却容颜不老，传说与她睡的寒玉床和修炼的古墓派内功有关。\par}

\plotframe{assets/levels/zhouzhiruo/level_16/frame_02.jpg}
  {周芷若线第十六拍《终南古墓》第2帧}
  {fig:level:zhouzhiruo:16:2}
  {\scriptsize}
  {\textbf{隐士：}周掌门，世人皆道古墓派驻颜有术——其实不在玉石，而在心静。小龙女之所以容颜不老，是因她心如止水、不染尘俗。女子养颜，五分在养身，五分在养心。心若焦躁，面便枯槁。\par}

\plotframe{assets/levels/zhouzhiruo/level_16/frame_03.jpg}
  {周芷若线第十六拍《终南古墓》第3帧}
  {fig:level:zhouzhiruo:16:3}
  {\scriptsize}
  {\mbox{}}

\clearpage

\plotframe{assets/levels/zhouzhiruo/level_16/frame_04.jpg}
  {周芷若线第十六拍《终南古墓》第4帧}
  {fig:level:zhouzhiruo:16:4}
  {\scriptsize}
  {\mbox{}}

\plotframe{assets/levels/zhouzhiruo/level_16/frame_05.jpg}
  {周芷若线第十六拍《终南古墓》第5帧}
  {fig:level:zhouzhiruo:16:5}
  {\scriptsize}
  {\mbox{}}

\subsubsection{第十七拍·华山论剑}\label{plot:zhouzhiruo:17}
本关以选择判断完成叙事向健康知识的转场，题目见\hyperref[question:zhouzhiruo:17]{第\ref*{question:zhouzhiruo:17}项}。

\plotframe{assets/levels/zhouzhiruo/level_17/frame_01.jpg}
  {周芷若线第十七拍《华山论剑》第1帧}
  {fig:level:zhouzhiruo:17:1}
  {\scriptsize}
  {\textbf{旁白：}华山论剑——武林中最负盛名的盛事。周芷若在此遇见来自各地的女侠，有人谈剑，有人论道，也有人说起了那个不为人知的话题。\par
\textbf{女侠乙：}唉，此等妇人隐秘之疾，江湖上少有人敢公开谈论。可不说——便是害人！\par}

\clearpage

\plotframe{assets/levels/zhouzhiruo/level_17/frame_02.jpg}
  {周芷若线第十七拍《华山论剑》第2帧}
  {fig:level:zhouzhiruo:17:2}
  {\scriptsize}
  {\textbf{女侠甲：}周掌门，你可知去年我师妹死于——宫颈癌？起因不过是一种极易感染的病毒（HPV），若是早些接种了疫苗，她本可活下去的……\par}

\plotframe{assets/levels/zhouzhiruo/level_17/frame_03.jpg}
  {周芷若线第十七拍《华山论剑》第3帧}
  {fig:level:zhouzhiruo:17:3}
  {\scriptsize}
  {\mbox{}}

\plotframe{assets/levels/zhouzhiruo/level_17/frame_04.jpg}
  {周芷若线第十七拍《华山论剑》第4帧}
  {fig:level:zhouzhiruo:17:4}
  {\scriptsize}
  {\mbox{}}

\clearpage

\plotframe{assets/levels/zhouzhiruo/level_17/frame_05.jpg}
  {周芷若线第十七拍《华山论剑》第5帧}
  {fig:level:zhouzhiruo:17:5}
  {\scriptsize}
  {\mbox{}}

\subsubsection{第十八拍·白莲重生}\label{plot:zhouzhiruo:18}
本关以选择判断完成叙事向健康知识的转场，题目见\hyperref[question:zhouzhiruo:18]{第\ref*{question:zhouzhiruo:18}项}。

\plotframe{assets/levels/zhouzhiruo/level_18/frame_01.jpg}
  {周芷若线第十八拍《白莲重生》第1帧}
  {fig:level:zhouzhiruo:18:1}
  {\scriptsize}
  {\textbf{旁白：}十八道关卡，十八次蜕变。周芷若走完了从汉水渔女→峨眉掌门→歧路失途→最终觉醒的全部旅程。她不再是谁的未婚妻、谁的徒弟、谁的对手——她只是她自己。\par
\textbf{旁白：}峨眉山的钟声响起——悠远、清澈、包容。这钟声曾经是复仇的战鼓，如今是祝福的梵音。\par}

\plotframe{assets/levels/zhouzhiruo/level_18/frame_02.jpg}
  {周芷若线第十八拍《白莲重生》第2帧}
  {fig:level:zhouzhiruo:18:2}
  {\scriptsize}
  {\textbf{周芷若：}（立于峨眉金顶，望着云海日出，微笑道）我这一生，最幸运的——不是遇见张无忌，而是在失去一切之后，找回了自己。从今往后，峨眉山下将开设义诊——愿天下女子，不必如我一般，在黑暗中独自摸索。\par}

\clearpage

\plotframe{assets/levels/zhouzhiruo/level_18/frame_03.jpg}
  {周芷若线第十八拍《白莲重生》第3帧}
  {fig:level:zhouzhiruo:18:3}
  {\scriptsize}
  {\mbox{}}

\plotframe{assets/levels/zhouzhiruo/level_18/frame_04.jpg}
  {周芷若线第十八拍《白莲重生》第4帧}
  {fig:level:zhouzhiruo:18:4}
  {\scriptsize}
  {\mbox{}}

\plotframe{assets/levels/zhouzhiruo/level_18/frame_05.jpg}
  {周芷若线第十八拍《白莲重生》第5帧}
  {fig:level:zhouzhiruo:18:5}
  {\scriptsize}
  {\mbox{}}

\clearpage
\subsubsection{路线结局：第十八拍 · 白莲重生}
\textbf{结尾诗}\begin{quote}汉水悠悠去，白莲岁岁新。\par 峨眉金顶上，云开见月明。\par 曾为情所困，终得心安宁。\par 愿世间女子，不困于深情。\end{quote}
\textbf{结局文本}\begin{quote}周芷若回到了峨眉山。在金顶之上，她望着云海翻涌、日出东方，终于放下了一切的执念。\par \par 她曾是汉水渔船上那个失去父亲的小女孩，曾是灭绝师太座下最听话的弟子，曾是那个在婚礼上被抛弃的新娘，曾是那个练九阴白骨爪、走上歧途的峨眉掌门。\par \par 如今，她只是周芷若。\par \par 她开始著书立说，将峨眉派医学典籍整理成册，开设义诊，为四方女子诊治疾病。\par \par 峨眉金顶的钟声日日响起——不是为了复仇，而是为了祝祷天下女子平安喜乐。\par \par 蔡文姬以《胡笳十八拍》叹乱世流离——周芷若以十八道关卡走过自己的一生，最终明白：真正的归汉，不是回到汉水，而是回到自己的内心。\end{quote}
